\documentclass[12pt]{article}
\usepackage[utf8]{inputenc}
\usepackage[T1]{fontenc}
\usepackage{hyperref}
\title{Pumpokol Language Resources for K-12 Curriculum Development}
\author{Compiled by Claude Sonnet 5.5 (claude-sonnet-5-5) for Andr\'as Kornai}
\date{}
\begin{document}
\maketitle

\section{Introduction}
Pumpokol (ISO 639-3 xpm; Glottocode pump1237) is an extinct Yeniseian language of Siberia. Sources I found give extinction in the early 18th century (about 1740 or earlier). It is poorly documented. Some material labelled Pumpokol has been identified as Yugh or Ketic, and one source reports O.\ Tailleur's view that it is a dialect of Ket. That view is not settled. No community of speakers, no school programme and no teaching materials turned up in my searches. A curriculum can only be a heritage-awareness unit built on the small attested word list.

\section{Primary Educational Resources}
I found no primers, textbooks or school programmes for Pumpokol. Nothing in the search results was designed for classroom use.

\section{Linguistic and Reference Materials}
\begin{itemize}
\item Wikipedia, ``Pumpokol language'': overview of status, classification and the small attested vocabulary (for example \emph{ab} `father', \emph{am} `mother'), drawing on Werner (2005). \url{https://en.wikipedia.org/wiki/Pumpokol_language}
\item Wikipedia, ``Pumpokolic languages'': the Pumpokolic group and its classification debate. \url{https://en.wikipedia.org/wiki/Pumpokolic_languages}
\item Wikipedia, ``Yeniseian languages'': the family context, including the Ket language, the only surviving Yeniseian language. \url{https://en.wikipedia.org/wiki/Yeniseian_languages}
\item ASJP database entry for Pumpokol: a short comparative word list. \url{https://asjp.clld.org/languages/PUMPOKOL}
\item An academia.edu item surfaced in the search results; I did not check its content. \url{https://www.academia.edu/86720593}
\end{itemize}

\section{Religious and Community Resources}
None found. No Pumpokol-specific community organisation or scripture translation appeared in the results.

\section{Digital and Audio Resources}
No audio or video recordings exist, since the language went extinct before recording technology. The ASJP entry above is the only digital word list found.

\section{Curriculum Development Context}
The evidence is too thin for a language-immersion curriculum. A realistic option is a short heritage unit on Yeniseian languages and the history of Siberian peoples. It could use the attested words from Werner (2005), as presented on the Wikipedia page, and be explicit about the uncertainty over what counts as Pumpokol. Any fuller work would need the primary sources and specialist Yeniseian linguists. I did not find contact details for any, and none are given here.
\end{document}